% =====================================================================
\chapter{Continuidad}
\label{tema:3}
% =====================================================================

Una función es continua cuando pequeños cambios en la variable producen pequeños cambios en el valor. Este tema formaliza esa idea con la definición de límite, demuestra que las operaciones usuales conservan la continuidad y prueba los resultados que Spivak agrupa bajo el título «tres teoremas difíciles»: el teorema de Bolzano (con su consecuencia, el de los valores intermedios) y el teorema de Weierstrass, que reúne la acotación de una función continua en un intervalo cerrado y la existencia de máximo y mínimo. Todos dependen de la completitud de los números reales. Se siguen los capítulos~6 a~8 de Spivak~\cite{spivak}, el capítulo~3 de Apostol~\cite{apostol}, el capítulo~4 de los apuntes de \emph{Analysis~I}~\cite{cambA1} y la §2.4 de OpenStax~\cite{osc1}.

\section{Funciones continuas}

\begin{definicion}[Continuidad]\label{def:continuidad}
Sea $f\colon D\to\R$ y $a\in D$. Se dice que $f$ es \emph{continua en $a$} si
\[
\forall\varepsilon>0\ \ \exists\delta>0\ \ \forall x\in D:\quad \abs{x-a}<\delta\ \Rightarrow\ \abs{f(x)-f(a)}<\varepsilon.
\]
Si $a$ es interior a $D$, esto equivale a $\lim_{x\to a}f(x)=f(a)$. Se dice que $f$ es \emph{continua} (en $D$) si lo es en cada punto de $D$.
\end{definicion}

La definición exige tres cosas a la vez: que $f(a)$ exista, que exista el límite y que coincidan. Por el \cref{teo:heine}, $f$ es continua en $a$ si y solo si $x_n\to a$ (con $x_n\in D$) implica $f(x_n)\to f(a)$: una función continua «conmuta con los límites», $\lim f(x_n)=f(\lim x_n)$.

\begin{proposicion}[Operaciones con funciones continuas]\label{prop:operaciones-continuas}
Si $f$ y $g$ son continuas en $a$, también lo son $f+g$, $cf$, $fg$ y, si $g(a)\ne0$, $f/g$. Si $f$ es continua en $a$ y $g$ es continua en $f(a)$, entonces $g\circ f$ es continua en $a$.
\end{proposicion}

\begin{proposicion}[Continuidad de las funciones elementales]\label{prop:elementales-continuas}
Son continuas en todo su dominio: los polinomios y las funciones racionales; $\abs{x}$; $\sin$ y $\cos$; la exponencial y el logaritmo; y, en consecuencia, $\tanh$, $\sigma$, $\ReLU$, $\softplus$ y todas las funciones que se obtienen de ellas por sumas, productos, cocientes y composiciones.
\end{proposicion}

\begin{definicion}[Tipos de discontinuidad]\label{def:discontinuidades}
Sea $f$ definida en un entorno de $a$ y no continua en $a$. La discontinuidad es
\begin{enumerate}
  \item \emph{evitable} si existe $\lim_{x\to a}f(x)$ y es finito (pero no coincide con $f(a)$, o $f(a)$ no está definido);
  \item \emph{de salto} si existen los límites laterales, finitos y distintos;
  \item \emph{esencial} en los demás casos (algún lateral infinito o inexistente).
\end{enumerate}
\end{definicion}

\addtocounter{definicion}{1}%

\begin{proposicion}[Conservación del signo]\label{prop:signo}
Si $f$ es continua en $a$ y $f(a)>0$, existe $\delta>0$ tal que $f(x)>0$ para todo $x\in D$ con $\abs{x-a}<\delta$.
\end{proposicion}

\section{El teorema de Bolzano y los valores intermedios}

\begin{teorema}[Bolzano]\label{teo:bolzano}
Sea $f\colon[a,b]\to\R$ continua con $f(a)<0<f(b)$. Entonces existe $c\in(a,b)$ con $f(c)=0$.
\end{teorema}

La demostración es constructiva: proporciona un algoritmo, la \emph{bisección}, que tras $n$ pasos localiza un cero con error menor que $(b-a)/2^{n}$.

\begin{corolario}[Teorema de los valores intermedios]\label{cor:tvi}
Sea $f\colon[a,b]\to\R$ continua y sea $y$ un valor estrictamente comprendido entre $f(a)$ y $f(b)$. Entonces existe $c\in(a,b)$ con $f(c)=y$. En consecuencia, la imagen de un intervalo por una función continua es un intervalo.
\end{corolario}

\addtocounter{definicion}{1}%

\addtocounter{definicion}{1}%

\begin{teorema}[Continuidad de la inversa]\label{teo:inversa-continua}
Sea $f\colon I\to\R$ continua y estrictamente monótona en un intervalo $I$. Entonces $J=f(I)$ es un intervalo y $f^{-1}\colon J\to I$ es continua.
\end{teorema}

Así, $\ln=\exp^{-1}$, $\sqrt[n]{\phantom{x}}$, $\arcsin$, $\arctan$ y $\logit=\sigma^{-1}$ son continuas.

\section{El teorema de Weierstrass}

\begin{lema}[Bolzano--Weierstrass]\label{lem:bw}
Toda sucesión acotada de números reales tiene una subsucesión convergente.
\end{lema}

\begin{teorema}[Weierstrass]\label{teo:weierstrass}
Sea $f\colon[a,b]\to\R$ continua. Entonces $f$ está acotada y alcanza sus valores máximo y mínimo: existen $x_m,x_M\in[a,b]$ con $f(x_m)\le f(x)\le f(x_M)$ para todo $x\in[a,b]$.
\end{teorema}

Las hipótesis son necesarias: en el intervalo abierto $(0,1)$, $f(x)=1/x$ es continua y no acotada, y $f(x)=x$ es acotada pero no alcanza ni su supremo ni su ínfimo; en $[0,1]$, la función que vale $1/x$ para $x>0$ y $0$ en~$0$ no es continua y no está acotada.

\addtocounter{definicion}{1}%

\addtocounter{definicion}{1}%

